\chapter{QED Renormalization and Electromagnetic Form Factors}
\label{ch:qed-renormalization-form-factors}
\ConventionsLorentzian


\section{Bare Variables and Ward Organization}

Use the electron convention of Chapter~\ref{ch:qed-external-states}. The bare
representative fields and
parameters are
\[
  A_\mu^B,\qquad \psi^B,\qquad \bar\psi^B,\qquad m^B,\qquad e^B,
\]
and the bare Lagrangian is
\[
  \mathcal L_B
  =
  -\frac14F^B_{\mu\nu}F_B^{\mu\nu}
  -
  \bar\psi^B
  \left[
    \gamma^\mu(\partial_\mu+\ii e^B A_\mu^B)+m^B
  \right]\psi^B .
\]
This is the explicit-coupling QED normalization: the photon representative has
a canonical kinetic term, and the electric charge appears in the covariant
derivative.  It differs from the coupling-absorbed Yang--Mills normalization of
Chapter~\ref{ch:vol2-classical-yang-mills-matter}.  In the notation of
\eqref{eq:ym-qed-curvature-rescaling}, the electron convention \(g=-e\) gives
\(A_\mu^{\rm YM}=-eA_\mu^{\rm QED}\) and
\(F_{\mu\nu}^{\rm YM}=-eF_{\mu\nu}^{\rm QED}\).  Therefore the photon
renormalization constant \(Z_A\) below is a field-residue constant in
explicit-coupling coordinates; after the rescaling to Yang--Mills coordinates
the same running is read from the coefficient of
\((F_{\mu\nu}^{\rm YM})^2/e^2\).
Introduce renormalized representative fields by
\[
  A_\mu^B=Z_A^{1/2}A_\mu,\qquad
  \psi^B=Z_\psi^{1/2}\psi,\qquad
  \bar\psi^B=Z_\psi^{1/2}\bar\psi .
\]
The constants \(Z_A\) and \(Z_\psi\) are fixed by pole-residue
normalizations in the infrared-regulated hard theory: \(Z_A\) normalizes the
one-photon component of \(\widehat F_{\mu\nu}\ket{\vac}\), and \(Z_\psi\)
normalizes the one-electron and one-positron spinor poles before the
long-distance regulator is removed.  In exact massless QED, \(Z_\psi\) is not
a physical residue of an isolated electron Wigner particle; its infrared
dependence is part of the hard-kernel bookkeeping that is combined with soft
radiation in Chapter~\ref{ch:volume-ii-infrared-inclusive-qed}. Write
\[
  m^B=m+\delta m .
\]

The renormalized charge parameter \(e\) is fixed by the electromagnetic
current matrix element at zero momentum transfer. Gauge covariance of the
renormalized representative derivative requires
\[
  e=e^B Z_A^{1/2},
\]
so that
\[
  \partial_\mu+\ii e^B A_\mu^B
  =
  \partial_\mu+\ii e A_\mu .
\]
Equivalently, if a separate vertex constant \(Z_1\) is introduced by
\[
  e^B Z_A^{1/2}Z_\psi=eZ_1,
\]
then Abelian gauge covariance gives the Ward organization
\[
  Z_1=Z_\psi .
\]
The same statement can be read from the renormalized representative gauge
transformation,
\[
  A_\mu\mapsto A_\mu+\partial_\mu\xi,\qquad
  \psi\mapsto \ee^{-\ii e\xi}\psi,\qquad
  \bar\psi\mapsto \bar\psi\,\ee^{\ii e\xi}.
\]
The coefficient \(e\), not \(eZ_\psi\), is the charge appearing in this
transformation law.  The common factor \(Z_\psi\) multiplies the whole spinor
covariant derivative and therefore belongs to the spinor residue
normalization.
Thus
\[
  \mathcal L_B
  =
  -\frac14Z_A F_{\mu\nu}F^{\mu\nu}
  -
  Z_\psi\bar\psi
  \left[
    \gamma^\mu(\partial_\mu+\ii e A_\mu)+m+\delta m
  \right]\psi .
\]
Split this into a renormalized quadratic and interaction part plus
counterterms:
\[
\begin{aligned}
  \mathcal L_B
  &=
  \underbrace{
    -\frac14F_{\mu\nu}F^{\mu\nu}
    -\bar\psi(\gamma^\mu\partial_\mu+m)\psi
    -\ii eA_\mu\bar\psi\gamma^\mu\psi
  }_{\mathcal L_{\mathrm{ren}}}
\\
  &\quad
  -\frac14(Z_A-1)F_{\mu\nu}F^{\mu\nu}
  -(Z_\psi-1)\bar\psi(\gamma^\mu\partial_\mu+m)\psi
  -Z_\psi\delta m\,\bar\psi\psi
\\
  &\quad
  -\ii e(Z_\psi-1)A_\mu\bar\psi\gamma^\mu\psi .
\end{aligned}
\]
The last line is the vertex counterterm tied to the spinor wavefunction
counterterm by the Abelian Ward identity. The mass counterterm is independent
because \(m\bar\psi\psi\) is gauge-invariant.

The current normalized by the renormalized field strength is obtained from
the Euler--Lagrange equation for \(A_\mu\):
\[
  \partial_\nu F^{\mu\nu}
  =
  -\ii Z_A^{-1}eZ_\psi\,\bar\psi\gamma^\mu\psi
  \equiv j^\mu .
\]
This identity is an operator identity at separated insertions in the regulated
theory.  In time-ordered products, differentiating the ordering symbols may
also produce local contact terms; the nonlocal momentum dependence discussed
below is insensitive to such polynomial contact terms.

\section{Photon Self-Energy and Transversality}

Let
\[
  \bra{\vac}
  \operatorname T\{\widehat A_\mu(x)\widehat A_\nu(0)\}
  \ket{\vac}
  =
  \int\frac{\dd^4k}{(2\pi)^4}\,
  \ee^{\ii k\cdot x}\,\widetilde G_{\mu\nu}(k).
\]
The full photon representative two-point function is obtained by inserting
one-particle-irreducible photon self-energies into the free propagator:
\[
\begin{gathered}
\begin{tikzpicture}[baseline=(current bounding box.center),>=Stealth,scale=0.88]
  \tikzset{
    photon/.style={decorate,decoration={snake,amplitude=1.15pt,segment length=5pt},line width=0.55pt}
  }
  \draw[photon] (-0.2,0)--(1.0,0);
  \node at (1.45,0) {\(=\)};
  \draw[photon] (1.9,0)--(3.1,0);
  \node at (3.55,0) {\(+\)};
  \draw[photon] (4.05,0)--(4.75,0);
  \node[circle,draw,minimum size=0.48cm,inner sep=0pt] at (5.05,0) {\scriptsize 1PI};
  \draw[photon] (5.35,0)--(6.05,0);
  \node at (6.5,0) {\(+\)};
  \draw[photon] (7.0,0)--(7.55,0);
  \node[circle,draw,minimum size=0.45cm,inner sep=0pt] at (7.82,0) {\scriptsize 1PI};
  \draw[photon] (8.1,0)--(8.65,0);
  \node[circle,draw,minimum size=0.45cm,inner sep=0pt] at (8.92,0) {\scriptsize 1PI};
  \draw[photon] (9.2,0)--(9.75,0);
  \node at (10.25,0) {\(\cdots\)};
\end{tikzpicture}
\end{gathered}
\]
Let \(\ii\Pi^{\mu\nu}(k)\) denote the sum of amputated one-particle
irreducible photon two-point insertions.  Current conservation implies
transversality in a way that does not depend on a gauge-dependent
longitudinal representative.  The Maxwell equation with the renormalized
current has the form
\[
  \partial_\nu \widehat F^{\mu\nu}=\widehat j^\mu,
  \qquad
  \partial_\mu\widehat j^\mu=0
\]
at separated points.  Let \(M^{\mu\nu}(k)\) be the noncontact coefficient in
the Fourier transform of
\(\bra{\vac}\operatorname T\widehat j^\mu(x)\widehat j^\nu(0)\ket{\vac}\).
Then
\[
  k_\mu M^{\mu\nu}(k)=0,\qquad k_\nu M^{\mu\nu}(k)=0 .
\]
On the other hand, applying
\(\partial_\alpha F^{\mu\alpha}\) to the photon Dyson series acts on each
external photon line by the transverse differential operator
\[
  k^2\eta^{\mu\rho}-k^\mu k^\rho .
\]
The local contact terms are polynomials in \(k\).  The nonlocal part is
therefore transverse only if the amputated 1PI insertion obeys
\[
  k_\mu\Pi^{\mu\nu}(k)=0 .
\]
Lorentz covariance then gives the tensor decomposition
\[
  \Pi^{\mu\nu}(k)
  =
  (\eta^{\mu\nu}k^2-k^\mu k^\nu)\pi(k^2)
\]
for a scalar function \(\pi\).

In covariant gauge, with gauge parameter \(\xi_{\mathrm g}\), the resummed
propagator is
\[
  \widetilde G_{\mu\nu}(k)
  =
  \frac{-\ii}{k^2-\ii\epsilon}
  \left[
    \frac{\eta_{\mu\nu}}{1-\pi(k^2)}
    -
    \frac{k_\mu k_\nu}{k^2-\ii\epsilon}
    \left(
      1-\xi_{\mathrm g}
      +\frac{\pi(k^2)}{1-\pi(k^2)}
    \right)
  \right].
\]
The longitudinal part depends on the representative gauge choice. The
transverse pole is normalized by imposing
\[
  \pi(0)=0
\]
for the massive-electron theory, where the loop function is regular at the
photon point \(k^2=0\).

\section{One-Loop Vacuum Polarization}

The photon counterterm contributes
\[
\begin{tikzpicture}[baseline=(current bounding box.center),>=Stealth,scale=0.82]
  \tikzset{photon/.style={decorate,decoration={snake,amplitude=1.1pt,segment length=5pt},line width=0.55pt}}
  \draw[photon] (-1.25,0)--(1.25,0);
  \draw[thick] (0,-0.18)--(0,0.18);
  \node[below] at (-1.25,-0.08) {\scriptsize \(\mu\)};
  \node[below] at (1.25,-0.08) {\scriptsize \(\nu\)};
\end{tikzpicture}
\quad
=
-\ii(Z_A-1)(k^2\eta_{\mu\nu}-k_\mu k_\nu),
\]
which contributes \(1-Z_A\) to \(\pi(k^2)\). The one-loop electron diagram is
\[
\begin{tikzpicture}[baseline=(current bounding box.center),>=Stealth,scale=0.9]
  \tikzset{
    photon/.style={decorate,decoration={snake,amplitude=1.1pt,segment length=5pt},line width=0.55pt},
    fermion/.style={postaction={decorate,decoration={markings,mark=at position .58 with {\arrow{Stealth}}}},line width=0.55pt}
  }
  \draw[photon] (-1.7,0)--(-0.62,0);
  \draw[fermion] (0,0) circle (0.62);
  \draw[photon] (0.62,0)--(1.7,0);
  \node[below] at (-1.7,-0.08) {\scriptsize \(\mu\)};
  \node[below] at (1.7,-0.08) {\scriptsize \(\nu\)};
  \node[above] at (0,0.70) {\scriptsize \(p\)};
  \node[below] at (0,-0.72) {\scriptsize \(p-k\)};
\end{tikzpicture}
\]
and equals
\[
\begin{aligned}
  \ii\Pi^{(1)}_{\mu\nu}(k)
  &=
  -
  \int\frac{\dd^4p}{(2\pi)^4}
  \operatorname{Tr}\left[
    e\gamma_\mu\,
    \frac{-\ii(-\ii\not p+m)}{p^2+m^2-\ii\epsilon}\,
    e\gamma_\nu\,
    \frac{-\ii(-\ii(\not p-\not k)+m)}
         {(p-k)^2+m^2-\ii\epsilon}
  \right].
\end{aligned}
\]
The minus sign is the closed fermion loop. Continue the integral to
\(D=4-\epsilon\) dimensions, keeping
\[
  \{\gamma_\mu,\gamma_\nu\}=2\eta_{\mu\nu},\qquad
  \eta^{\mu\nu}\eta_{\mu\nu}=D,\qquad
  \operatorname{Tr}\mathbf 1=4 .
\]
After Wick rotation and the Feynman-parameter identity
\[
  \frac1{ab}=\int_0^1\dd x\,
  \frac1{\bigl(xa+(1-x)b\bigr)^2},
\]
shift
\[
  p\mapsto p+xk .
\]
The denominator becomes
\[
  \bigl(p^2+x(1-x)k^2+m^2\bigr)^2 .
\]
The trace identity
\[
  \operatorname{Tr}(
    \gamma_\mu\gamma_\alpha\gamma_\nu\gamma_\beta)
  =
  4(
    \eta_{\mu\alpha}\eta_{\nu\beta}
    -\eta_{\mu\nu}\eta_{\alpha\beta}
    +\eta_{\mu\beta}\eta_{\nu\alpha})
\]
reduces the numerator. Rotational symmetry in the continued loop momentum
allows
\[
  p_\mu p_\nu\mapsto \frac1D\eta_{\mu\nu}p^2
\]
inside the integral, and all terms linear in \(p\) vanish.  With
\[
  M_x^2=m^2+x(1-x)k^2,
\]
the angularly averaged numerator can be written as
\[
\begin{aligned}
  N_{\mu\nu}
  =
  4\Bigl[
    &\eta_{\mu\nu}\bigl(p^2-x(1-x)k^2+m^2\bigr)
    -\frac2D\eta_{\mu\nu}p^2
\\
    &\qquad\qquad
    +2x(1-x)k_\mu k_\nu
  \Bigr].
\end{aligned}
\]
The two scalar integrals needed are
\[
\begin{aligned}
  \int\frac{\dd^Dp}{(2\pi)^D}\frac1{(p^2+M_x^2)^2}
  &=
  \frac1{(4\pi)^{D/2}}
  \Gamma(2-D/2)(M_x^2)^{D/2-2},\\
  \int\frac{\dd^Dp}{(2\pi)^D}\frac{p^2}{(p^2+M_x^2)^2}
  &=
  \frac{D}{2}\frac1{(4\pi)^{D/2}}
  \Gamma(1-D/2)(M_x^2)^{D/2-1},
\end{aligned}
\]
where the second formula is defined by analytic continuation in \(D\).
Substitution cancels the nontransverse local-looking pieces and leaves the
transverse tensor
\[
  \Pi^{(1)}_{\mu\nu}(k)
  =
  -(\eta_{\mu\nu}k^2-k_\mu k_\nu)
  \frac{8e^2}{(4\pi)^{D/2}}
  \Gamma(2-D/2)
  \int_0^1\dd x\,x(1-x)
  \bigl(m^2+x(1-x)k^2\bigr)^{D/2-2}.
\]
With the counterterm included,
\[
  \pi(k^2)
  =
  1-Z_A
  -
  \frac{8e^2}{(4\pi)^{D/2}}\Gamma(2-D/2)
  \int_0^1\dd x\,x(1-x)
  \bigl(m^2+x(1-x)k^2\bigr)^{D/2-2}.
\]
The pole normalization \(\pi(0)=0\) fixes
\[
  Z_A
  =
  1-
  \frac{8e^2}{(4\pi)^{D/2}}\Gamma(2-D/2)\frac16\,m^{D-4}.
\]
For \(D=4-\epsilon\), this has the pole part
\[
  Z_A=1-\frac{e^2}{6\pi^2}\frac1\epsilon+\text{finite}.
\]
Subtracting the value at \(k^2=0\) and taking \(D\to4\) gives the finite
one-loop vacuum-polarization function
\[
  \pi(k^2)
  =
  \frac{e^2}{2\pi^2}
  \int_0^1\dd x\,x(1-x)
  \log
  \frac{m^2+x(1-x)k^2}{m^2}.
\]
This \(\pi(k^2)\) is the cutoff-removed form factor: in a cutoff
realization it is the \(\Lambda\to\infty\) limit of the renormalized
transverse coefficient in the photon two-point matrix element, with the
charge normalization \(\pi(0)=0\) held fixed.  In dimensional
regularization the same statement is represented by the analytic continuation
to \(D=4\) after the local counterterm \(Z_A\) has subtracted the ultraviolet
pole.

\section{The Electron Current Matrix Element}

Let \(|\vec p,\sigma\rangle\) be a one-electron state of charge \(q=-e\).
Translation covariance gives
\[
  \langle\vec p',\sigma'|\widehat j^\mu(x)|\vec p,\sigma\rangle
  =
  \ee^{\ii(p-p')\cdot x}
  \langle\vec p',\sigma'|\widehat j^\mu(0)|\vec p,\sigma\rangle .
\]
Write the matrix element as
\[
  \langle\vec p',\sigma'|\widehat j^\mu(0)|\vec p,\sigma\rangle
  =
  \frac{\ii q}{(2\pi)^3}
  \bar u^{\sigma'}(\vec p')M^\mu(p',p)u^\sigma(\vec p).
\]
Modulo the on-shell Dirac equations
\[
  (\not p-\ii m)u^\sigma(\vec p)=0,\qquad
  \bar u^{\sigma'}(\vec p')(\not p'-\ii m)=0,
\]
a parity-even Lorentz vector between spinor solutions can be reduced, by the
Clifford algebra, to the span of
\[
  \gamma^\mu,\qquad (p+p')^\mu,\qquad (p-p')^\mu .
\]
Terms with \(\gamma_5\) are excluded by parity, and terms with
\([\gamma^\mu,\gamma^\nu]\) are equivalent to the displayed basis by the
Gordon identity.  Current conservation removes the component proportional to
\((p-p')^\mu\). Therefore
\[
  M^\mu(p',p)
  =
  \gamma^\mu F(k^2)
  -\frac{\ii}{2m}(p+p')^\mu G(k^2),
  \qquad
  k=p-p' .
\]
The scalar functions \(F\) and \(G\) are the electromagnetic form factors in
the present gamma-matrix convention and in the \((p+p')^\mu\) basis.  For
comparison with the Dirac--Pauli basis, let
\(F_1^{\rm DP}\) and \(F_2^{\rm DP}\) be defined in this same gamma-matrix
convention by
\[
  \bar u(p')M^\mu u(p)
  =
  \bar u(p')
  \left[
    \gamma^\mu F_1^{\rm DP}(k^2)
    +
    \left(\gamma^\mu+\frac{\ii}{2m}(p+p')^\mu\right)
    F_2^{\rm DP}(k^2)
  \right]
  u(p).
\]
The second structure is the Gordon-identity form of the Pauli tensor
structure; its matrix element vanishes at \(p'=p\).  Comparing coefficients
with the displayed \((F,G)\) decomposition gives
\[
  F_1^{\rm DP}=F+G,
  \qquad
  F_2^{\rm DP}=-G,
  \qquad
  F=F_1^{\rm DP}+F_2^{\rm DP}.
\]
Thus the charge normalization below, \(F(0)+G(0)=1\), is precisely
\(F_1^{\rm DP}(0)=1\), while the anomalous magnetic moment is
\(F_2^{\rm DP}(0)=-G(0)\).

Charge normalization fixes one combination at zero momentum transfer. Taking
\(p'\to p\), using
\[
  \bar u^{\sigma'}(\vec p)u^\sigma(\vec p)
  =
  \frac{m}{p^0}\delta_{\sigma\sigma'},
\]
and using the Dirac equations to rewrite
\[
  p^\mu\bar u^{\sigma'}(\vec p)u^\sigma(\vec p)
  =
  \ii m\,\bar u^{\sigma'}(\vec p)\gamma^\mu u^\sigma(\vec p),
\]
one obtains
\[
  \langle\vec p,\sigma'|\widehat j^\mu(0)|\vec p,\sigma\rangle
  =
  \frac{q}{(2\pi)^3}
  \frac{p^\mu}{p^0}
  \delta_{\sigma\sigma'}\bigl(F(0)+G(0)\bigr).
\]
The charge operator
\[
  \widehat Q=\int\dd^3\vec x\,\widehat j^0(x)
\]
acts by
\[
  \widehat Q|\vec p,\sigma\rangle=q|\vec p,\sigma\rangle .
\]
Thus
\[
  F(0)+G(0)=1.
\]
The equality is a distributional momentum-state identity: after smearing the
external plane waves by a common wave packet and then removing the smearing,
the spatial integral supplies the same
\(\delta^{(3)}(\vec p'-\vec p)\) on both sides.

\section{Magnetic Moment from the Form Factor}

Couple the current to a weak classical background:
\[
  \Delta H
  =
  -\int\dd^3\vec x\,A_\mu^{\mathrm{cl}}(x)\widehat j^\mu(x).
\]
For a slowly varying magnetic field choose
\[
  A_0^{\mathrm{cl}}=0,\qquad
  \nabla\times\vec A^{\mathrm{cl}}=\vec B .
\]
The part of the form-factor matrix element proportional to
\((p+p')^\mu\) gives the spin-independent Lorentz-force contribution in the
nonrelativistic limit.  The spin-dependent term is exposed by writing the
\(\gamma^\mu F\) term with the Gordon identity.  Equivalently, after an
integration by parts,
\[
  \int\dd^3\vec x\,
  A_i^{\mathrm{cl}}(x)\,
  \ee^{\ii(\vec p-\vec p')\cdot\vec x}(p'-p)_jS^{ij}
  =
  \frac{\ii}{2}\int\dd^3\vec x\,
  \ee^{\ii(\vec p-\vec p')\cdot\vec x}
  F^{\mathrm{cl}}_{ij}(x)S^{ij},
\]
up to terms higher order in the spatial momentum transfer; only the
antisymmetric part contributes to the spin coupling.
The spin-dependent part of
\(\langle\vec p',\sigma'|\Delta H|\vec p,\sigma\rangle\) is obtained from
the identity
\[
\begin{aligned}
  &\bar u^{\sigma'}(\vec p')
  [\gamma^\mu,\gamma^\nu](p'-p)_\nu
  u^\sigma(\vec p)  \\
  &\quad =
  -4\ii m\,\bar u^{\sigma'}(\vec p')\gamma^\mu u^\sigma(\vec p)
  +2(p^\mu+p'^\mu)
  \bar u^{\sigma'}(\vec p')u^\sigma(\vec p),
\end{aligned}
\]
with
\[
  S^{\mu\nu}=-\frac{\ii}{4}[\gamma^\mu,\gamma^\nu].
\]
In the nonrelativistic limit \(|\vec p|,|\vec p'|\ll m\), the spin
coupling becomes
\[
  \Delta H_{\mathrm{spin}}
  =
  -\frac{qF(0)}{2m}\,\vec B\cdot\vec\sigma .
\]
Writing this as
\[
  \Delta H_{\mathrm{spin}}
  =
  -g_{\mathrm{mag}}\frac{q}{2m}\,
  \vec B\cdot\frac{\vec\sigma}{2},
\]
the magnetic factor is
\[
  g_{\mathrm{mag}}=2F(0).
\]
Together with \(F(0)+G(0)=1\), this gives
\[
  g_{\mathrm{mag}}=2\bigl(1-G(0)\bigr).
\]

\section{One-Loop Vertex Correction}

The current is represented by
\[
  \widehat j^\mu
  =
  -\ii Z_A^{-1}e\,Z_\psi\,\bar\psi\gamma^\mu\psi
  =
  -\ii e^B Z_A^{-1/2}Z_\psi\,\bar\psi\gamma^\mu\psi .
\]
Within the same infrared-regulated hard-kernel convention, the current
insertion between external electron labels is computed by reducing the
external spinor legs by LSZ:
\[
\begin{tikzpicture}[baseline=(current bounding box.center),>=Stealth,scale=0.88]
  \tikzset{
    fermion/.style={postaction={decorate,decoration={markings,mark=at position .58 with {\arrow{Stealth}}}},line width=0.55pt},
    photon/.style={decorate,decoration={snake,amplitude=1.15pt,segment length=5pt},line width=0.55pt}
  }
  \coordinate (v) at (0,0);
  \draw[fermion] (v)--(1.0,0.65);
  \draw[fermion] (1.0,-0.65)--(v);
  \draw[photon] (-0.95,0)--(v);
  \draw[fill=white,line width=0.55pt] (v) circle (0.1);
  \draw[line width=0.55pt] (-0.07,-0.07)--(0.07,0.07);
  \draw[line width=0.55pt] (-0.07,0.07)--(0.07,-0.07);
  \node[left] at (-0.95,0) {\scriptsize \(j^\mu\)};
  \node[right] at (1.05,0.65) {\scriptsize \(\bar u(p')\)};
  \node[right] at (1.05,-0.65) {\scriptsize \(u(p)\)};
\end{tikzpicture}
\]
If the three-point function is written in bare fields, each external fermion
leg supplies a factor \(Z_\psi^{1/2}\) in the LSZ residue.  If it is written
in renormalized fields, those factors are absent from the external residues
but reappear through the counterterm vertices.  These two bookkeepings give
the same matrix element.
After removing the infrared regulator in massless QED, this exclusive
one-electron matrix element is not an autonomous physical scattering
amplitude.  The hard form factor is one factor in either an inclusive
probability, where unresolved photons are summed over, or a dressed-state
matrix element whose charged asymptotic state already contains the soft
electromagnetic field.

The tree current vertex gives
\[
  -\ii Z_A^{-1}e\,\frac{Z_\psi}{(2\pi)^3}
  \bar u^{\sigma'}(\vec p')\gamma^\mu u^\sigma(\vec p).
\]
The current insertion can also be followed by any number of photon
self-energy insertions before the photon attaches to the electron line.  The
transverse photon chain contributes \(Z_A/(1-\pi(k^2))\), so the whole
tree-plus-vacuum-polarization chain contributes
\[
  -\ii e\,\frac{Z_\psi}{(2\pi)^3}
  \frac1{1-\pi(k^2)}
  \bar u^{\sigma'}(\vec p')\gamma^\mu u^\sigma(\vec p).
\]
This changes only the coefficient of \(\gamma^\mu\), hence only \(F(k^2)\).
It does not contribute to \(G(k^2)\).

At order \(e^2\), the contribution to \(G(k^2)\) comes from the vertex
correction
\[
\begin{tikzpicture}[baseline=(current bounding box.center),>=Stealth,scale=0.95]
  \tikzset{
    fermion/.style={postaction={decorate,decoration={markings,mark=at position .58 with {\arrow{Stealth}}}},line width=0.55pt},
    photon/.style={decorate,decoration={snake,amplitude=1.15pt,segment length=5pt},line width=0.55pt}
  }
  \coordinate (v) at (0,0);
  \coordinate (a) at (0.9,0.72);
  \coordinate (b) at (0.9,-0.72);
  \draw[fermion] (v)--(a);
  \draw[fermion] (b)--(v);
  \draw[photon] (a)--(b);
  \draw[fermion] (a)--(1.9,1.08);
  \draw[fermion] (1.9,-1.08)--(b);
  \draw[photon] (-0.85,0)--(v);
  \draw[fill=white,line width=0.55pt] (v) circle (0.1);
  \draw[line width=0.55pt] (-0.07,-0.07)--(0.07,0.07);
  \draw[line width=0.55pt] (-0.07,0.07)--(0.07,-0.07);
  \node[left] at (-0.85,0) {\scriptsize \(k=p-p'\)};
  \node[above] at (1.45,0.98) {\scriptsize \(p'-q\)};
  \node[below] at (1.45,-0.98) {\scriptsize \(p-q\)};
  \node[right] at (1.05,0) {\scriptsize \(q\)};
\end{tikzpicture}
\]
in Feynman gauge. Its integrand is
\[
\begin{aligned}
  &-\ii Z_A^{-1}e\,\frac{Z_\psi}{(2\pi)^3}
  \int\frac{\dd^4q}{(2\pi)^4}
  \frac{-\ii}{q^2-\ii\epsilon}\,
  \bar u(p')\,e\gamma^\rho
  \frac{-\ii[-\ii(\not p'-\not q)+m]}
       {(p'-q)^2+m^2-\ii\epsilon}
\\
  &\hspace{9em}\times
  \gamma^\mu
  \frac{-\ii[-\ii(\not p-\not q)+m]}
       {(p-q)^2+m^2-\ii\epsilon}
  e\gamma_\rho\,u(p).
\end{aligned}
\]
To isolate the term that affects the magnetic moment, Wick rotate, continue
the loop integral dimensionally, combine the three denominators by Feynman
parameters, and use the external Dirac equations.  The reduced numerator has
the form
\[
  \bar u(p')\gamma^\mu u(p)\,A(k^2)
  +(p+p')^\mu\bar u(p')u(p)\,B(k^2).
\]
The scalar \(A(k^2)\) contains ultraviolet and infrared singular terms whose
contribution is fixed by the charge normalization and the spinor
wavefunction counterterm.  The scalar \(B(k^2)\) is finite at this order and
determines \(G(k^2)\).  The resulting coefficient is
\[
  G(k^2)
  =
  -\frac{e^2m^2}{4\pi^2}
  \int_0^1\dd x\int_0^x\dd y\,
  \frac{x(1-x)}
       {m^2x^2+k^2y(x-y)} .
\]
The displayed \(G(k^2)\) is again the cutoff-removed form factor: for fixed
external on-shell wave packets and fixed nonzero momentum-transfer invariant
away from infrared singular loci, it is the ultraviolet-regulator removal
limit of the renormalized coefficient multiplying
\((p+p')^\mu\bar u(p')u(p)\) in the current matrix element.  The charge
normalization and the wavefunction counterterm determine the accompanying
\(\gamma^\mu\) coefficient, while no additional ultraviolet subtraction is
needed for this one-loop \(G(k^2)\).
At zero momentum transfer,
\[
  \int_0^1\dd x\int_0^x\dd y\,
  \frac{x(1-x)}{m^2x^2}
  =
  \frac1{2m^2},
  \qquad
  G(0)
  =
  -\frac{e^2}{8\pi^2}.
\]
With
\[
  \alpha=\frac{e^2}{4\pi},
\]
the one-loop magnetic factor is
\[
  g_{\mathrm{mag}}
  =
  2\bigl(1-G(0)\bigr)
  =
  2+\frac{\alpha}{\pi}+O(\alpha^2).
\]
